\subsection{Compact-boson wavefunction and Haldane--Shastry checks}
\label{sec:boson-sanity}
We evaluate Eq.~\eqref{eq:bosonpolynomial} independently of the
logarithmic correlator implementation and compare normalized vectors.
For each fixed $N$, charge pair $(a,b)$, and endpoint $W$, let $\psi$
be the normalized vector obtained by directly evaluating the
Laughlin/Jastrow polynomial in Eq.~\eqref{eq:bosonpolynomial}
(\texttt{polynomial\_state} in the code). Let $\phi$ be the normalized
conformal-block vector obtained by exponentiating the logarithmic
amplitudes in Eq.~\eqref{eq:logblock} (\texttt{vertex\_state}). Both
vectors represent the same chosen primary, with identical coordinates,
endpoint and occupation basis. Define their phase-aligned amplitude error
\begin{equation}
 e_{\rm amp}(\psi,\phi)=
 \max_O\left|\psi(O)e^{-i\arg\langle\phi|\psi\rangle}-\phi(O)\right|.
 \label{eq:amplitudeerror}
\end{equation}
The argument $\arg$ is the phase of the complex overlap. This comparison
checks that two implementations produce the same wavefunction, including
its phases and normalization. To test an independent spectral property
of that wavefunction, we also examine its eigenstate residual.

\paragraph{Why use the Haldane--Shastry Hamiltonian?}
The choice is specific to the exponent-two Jastrow states on the uniform
circle used here. Their half-filled vacuum is exactly the ground state
of the Haldane--Shastry (HS) spin chain, whose continuum description is
the compact boson at the $SU(2)_1$ radius. The conformal-block vacuum
and its relation to HS are given in Ref.~\cite{Herwerth}, Sec.~II A,
Eq.~(10), and Sec.~III A, Eq.~(37). This makes HS a natural additional
benchmark for our three explicit lattice states. In our energy convention,
the Hamiltonian is
\begin{equation}
 H_{\rm HS}=\left(\frac{\pi}{N}\right)^2
 \sum_{j<k}\frac{\bm S_j\cdot\bm S_k}{\sin^2[\pi(j-k)/N]},
 \qquad \bm S_j=(X_j,Y_j,Z_j)/2.
 \label{eq:HHS}
\end{equation}
Here $X,Y,Z$ are Pauli matrices, so each site is a spin one-half.
For each normalized state we compute
\begin{equation}
 E_{\rm HS}=\langle\psi|H_{\rm HS}|\psi\rangle,\qquad
 \eta_{\rm HS}=\|(H_{\rm HS}-E_{\rm HS})\ket{\psi}\|_2.
\end{equation}
This is Eq.~\eqref{eq:eigenresidual} with the Hamiltonian specified.
A small $\eta_{\rm HS}$ tests stationarity under HS; it does not by
itself identify the CFT primary or establish its conformal weights.
For the three states at $N=8,10,12$, the largest residual at strict
infinity is $6.21\times10^{-15}$. The excited-state residuals at
$W=10^8$ range from $4.60\times10^{-9}$ to $6.30\times10^{-9}$,
showing the small departure from an HS eigenstate introduced by retaining
the original finite endpoint.

All three boson vectors are constructed directly from the conformal-block
or polynomial amplitudes before applying this test. The HS Hamiltonian
is used only as a diagnostic: it neither generates these vectors nor
enters the relative entropy, which is calculated from their reduced
density matrices. Its use here is not a requirement of the continuum
relative-entropy theory and does not extend automatically to other
compactification radii or lattice geometries.

\begin{table}[htbp]\centering\small
\begin{tabular}{rlrrrr}\toprule
$N$ & State & $e_{\rm amp}(\infty)$ & $e_{\rm amp}(10^8)$ & $\eta(\infty)$ & $\eta(10^8)$\\\midrule
8 & $00$ & $9.30\times10^{-16}$ & $9.30\times10^{-16}$ & $1.04\times10^{-15}$ & $1.04\times10^{-15}$ \\
8 & $ss$ & $1.04\times10^{-15}$ & $8.50\times10^{-16}$ & $1.43\times10^{-15}$ & $5.97\times10^{-9}$ \\
8 & $3s,s$ & $4.38\times10^{-16}$ & $5.00\times10^{-16}$ & $8.02\times10^{-16}$ & $6.30\times10^{-9}$ \\
10 & $00$ & $2.86\times10^{-15}$ & $2.86\times10^{-15}$ & $1.58\times10^{-15}$ & $1.58\times10^{-15}$ \\
10 & $ss$ & $1.18\times10^{-15}$ & $1.14\times10^{-15}$ & $1.32\times10^{-15}$ & $5.22\times10^{-9}$ \\
10 & $3s,s$ & $1.46\times10^{-15}$ & $2.01\times10^{-15}$ & $9.82\times10^{-16}$ & $5.35\times10^{-9}$ \\
12 & $00$ & $2.89\times10^{-15}$ & $2.89\times10^{-15}$ & $6.21\times10^{-15}$ & $6.21\times10^{-15}$ \\
12 & $ss$ & $1.70\times10^{-15}$ & $2.24\times10^{-15}$ & $5.58\times10^{-15}$ & $4.60\times10^{-9}$ \\
12 & $3s,s$ & $1.88\times10^{-15}$ & $1.83\times10^{-15}$ & $5.11\times10^{-15}$ & $4.66\times10^{-9}$ \\
\bottomrule\end{tabular}

\caption{Direct polynomial versus logarithmic conformal-block evaluation
for the three canonical boson states. Here $\psi$ is the normalized
direct Laughlin/Jastrow polynomial vector of Eq.~\eqref{eq:bosonpolynomial},
and $\phi$ is the normalized conformal-block vector obtained by
exponentiating Eq.~\eqref{eq:logblock}. Each comparison uses the same
$N$, primary and endpoint $W$ for both vectors. The two amplitude-error
columns give $e_{\rm amp}(\psi,\phi)$ from Eq.~\eqref{eq:amplitudeerror};
the two residual columns apply Eqs.~\eqref{eq:eigenresidual} and
\eqref{eq:HHS} to $\psi$.
All vectors are normalized, $z_j=e^{2\pi i j/N}$, and the endpoint is
either strict infinity or $10^8$. The finite-endpoint residuals quantify
the perturbation retained to match the original Laughlin implementation.
No interval length or mixing probability enters these whole-state tests.}
\label{tab:bosoncanonical}
\end{table}
